\chapter{硬件模块框图与接口实现}
\section{硬件平台及模块连接}
实物为 Digilent Nexys A7-100T，PCB 标识为 PB200-292 REV D，器件目标为 \code{xc7a100tcsg324-1}。本项目使用板载传感器和人机接口，以 USB 连接电脑；FPGA 内部重构采集、分析、记录与通信电路。官方板卡资料用于确认接口连接，最终约束文件定义实际管脚\cite{rm,xdc,sch}。

\fig{board}{开发板实物与主要功能区域}{照片由操作者提供；标注为区域导览，不据照片推断未使用跳线的电气用途。}

\fig{hardware}{硬件功能模块与接口连接框图}{箭头方向以数据/信号流为准。I²C为开漏双向线路；XADC为FPGA内部专用模块，不占用外部传感器管脚。}

板外连接集中在 PROG/UART 接口。FT2232 的配置通道用于加载位流，UART通道用于运行数据；USB HOST、VGA、以太网及耳机口不属于当前自研系统的数据通路。原官方出厂演示用于早期板卡基线检查，本系统的完整录音在电脑端保存和播放。

\section{关键管脚与电气连接}
\begin{table}[H]\centering\caption{采集与通信关键引脚（FPGA视角）}\footnotesize\renewcommand{\arraystretch}{1.12}
\begin{tabularx}{\linewidth}{p{29mm}p{30mm}p{12mm}Y}\toprule
模块/信号 & FPGA管脚 & 方向 & 连接和职责\\\midrule
CLK100MHZ & E3 & 输入 & 板载100 MHz时钟，建立10 ns周期约束。\\
M\_CLK / M\_DATA & J5 / H5 & 出/入 & 麦克风时钟与PDM比特流。\\
M\_LRSEL & F5 & 输出 & PDM通道选择固定为0。\\
ACL\_SCLK / CSN & F15 / D15 & 输出 & ADXL362 SPI时钟与低有效片选。\\
ACL\_MOSI / MISO & F14 / E15 & 出/入 & SPI命令发送与XYZ读回。\\
ACL\_INT1 & B13 & 输入 & 数据就绪中断，触发读取并记录时刻。\\
TMP\_SCL / SDA & C14 / C15 & 双向 & ADT7420 I²C开漏线路，7位地址0x4B。\\
TMP\_INT / CT & D13 / B14 & 输入 & 温度告警相关电平监测。\\
{\footnotesize UART\_TXD\_IN} & C4 & 输入 & FTDI发送至FPGA接收。\\
{\footnotesize UART\_RXD\_OUT} & D4 & 输出 & FPGA发送至FTDI接收。\\
CPU\_RESETN & C12 & 输入 & 低有效板级复位。\\
BTNU/D/L/R/C & M18 / P18 / P17 / M17 / N17 & 输入 & 五向按键，其中BTNC供收藏与游戏刹车使用。\\
SW15 & V10 & 输入 & 识曲模式快捷进入与退出。\\\bottomrule
\end{tabularx}\end{table}

表中接口使用 LVCMOS33；全部普通开关中，SW8/T8 与 SW9/U8 使用 LVCMOS18，必须遵循各自 Bank 电压。完整16路开关、16路LED、RGB和数码管的引脚表随 XDC 交付，报告保留关键连接，避免重复抄录大量同类管脚。

输入同步与去抖承担不同职责。两级同步器降低异步输入对逻辑域的影响；开关按键再经过稳定计数去抖，输出状态及按下/释放计数。PDM、SPI输入及外部中断使用相应同步边界，不能直接把外部引脚当作理想同步信号。数码管位选、段选和LED属于面向人的显示输出，与采集链路分离。

\section{音频采集：PDM到PCM}
数字麦克风输入是密度编码比特流，不能直接作为16位声音样本使用。系统由100 MHz时钟产生2 MHz麦克风时钟，采集的PDM位映射为$+1/-1$，再通过三级积分梳状滤波器（CIC）按125倍抽取：
\begin{equation}
 f_\mathrm{PCM}=\frac{2\times10^6}{125}=16000\ \mathrm{Hz}.
\end{equation}
该结构以加减与延时实现主要低通及抽取，适合FPGA流水处理。三级积分器采用24位模运算，抽取后的三级梳状差分恢复输出；定点缩放使用对称舍入并饱和到 PCM16。音频饱和次数单独计数，便于发现过大输入。

\fig{audio}{音频多速率处理与两类数据产品}{16 kHz原速路径供事件保存和FFT快照；31抽头FIR及二倍抽取只服务长时间音乐录音。}

采集链路同时生成预览包络和原始快照。预览每16个PCM点保存最小/最大值，形成1 ms一组的轻量曲线；频谱使用单独冻结的512点原速声音快照。两者职责不同：包络用于观察变化，FFT用于观察短时频率结构，完整录音用于播放与识别。

\section{加速度、温度和片内监测}
ADXL362由2 MHz、模式0的SPI控制器访问，配置$\pm2g$量程和标称200 Hz输出率\cite{adxl}。启动时先读取器件身份，再写配置并回读核对；数据就绪后连续读回XYZ六字节，并保存实际时刻。传感器具有独立内部时钟，实际间隔不强行替换成5 ms，因此姿态和声振配对能够基于真实数据时刻工作。

ADT7420通过I²C读写寄存器。控制器以开漏方式驱动SCL/SDA，只主动拉低线路，释放后由上拉形成高电平；接收端同步读取实际线路状态，处理ACK、重复START、时钟拉伸及超时。系统核对身份寄存器、设置16位温度模式并回读，正常时周期读数；通信异常进入恢复等待后重新识别\cite{adt}。温度换算为$T=\mathrm{raw}/128$。

XADC采用连续序列器读取片内温度、VCCINT、VCCAUX及VCCBRAM，配置16次平均\cite{xadc}。有效位与原始码一起传给主机，分别换算：
\begin{equation}
 T_\mathrm{die}=\frac{503.975\,r_T}{65536}-273.15,\qquad
 V=\frac{3r_V}{65536}.
\end{equation}
板载ADT温度与FPGA结温来自不同位置，应分别显示；ADXL内部温度保留原始计数。换算是使用器件标称关系，不替代外部仪表的绝对标定。

\section{FPGA内部通路与记录缓存}
顶层将接口控制器、特征核心、触发逻辑、记录存储、监测快照及串口命令连接起来。各模块共用100 MHz计数器，PCM样本以固定间隔推进，XYZ使用实到时间。控制路径负责参数、状态和存储生命周期，数据路径负责按时接收、计算和保存，避免界面或网络操作反向阻塞板端采样。

\fig{buffers}{事件双组缓存与音乐独立缓冲的组织}{环形保存最近数据；触发后补足后段并冻结。主机读完且校验成功后确认释放；音乐与游戏等应用通过模式管理协调。}

事件缓冲维护连续滚动的前0.5秒历史；触发时记录锚点，再补足1秒后段，形成1.5秒完整窗口。声音固定24000点，XYZ按窗口和实际时刻筛选。两组缓冲允许读取一个冻结事件时，另一组继续承担记录；容量耗尽时必须给出状态，不能把有限缓存解释为无限连续存储。

音乐使用最多160000点的独立缓冲。10秒模式保存16 kHz，20秒模式保存经FIR处理的8 kHz；冻结后可按编号和偏移反复读取。中止命令只针对当前录音编号处理，释放未完成任务；已经保存并校验的完整音频在电脑端继续保留。

\section{板载显示与时序实现}
板载反馈包括16路LED、两组RGB和八位数码管。显示控制器将任务阶段投影为 \code{tAP3}、\code{rEC}、\code{SEnd}、\code{FInd}、\code{PASS}等状态；录音计数取自板端，传输进度取自主机实际导出量，避免用固定动画替代真实完成度。位扫描包含消隐和互斥，PWM控制名义亮度；显示熄灭不停止采样。

数字电路以10 ns系统时钟约束实现；真正异步的输入只在同步器边界设置例外，慢速串行输出设置20 ns传播上限，面向人的显示输出按其异步使用性质处理。最终实现的建立与保持裕量均为正。具体资源、时序与测试口径见第5章，完整约束和实现报告随源码交付。
